\section{Necessary Conditions on Any Frontier Statement}\label{sec:natural-families}

The realizability barrier leaves a more precise frontier question. If quotient shape alone is too expressive, what should replace it as the organizing unit of the metatheorem?

At minimum, the tractability frontier should not classify arbitrary sets of decision problems. It should range over representation-level structural classes and it should ignore benign changes of presentation.

The right eventual condition list will almost certainly be stronger than this minimal principle. For example, a mature frontier theorem may also require closure under adding explicitly irrelevant coordinates, adjoining dummy actions, or other representation-preserving refinements. But action/state relabel invariance is the irreducible starting point.

\leanmetapending{\LH{FR30}, \LH{FR31}, \LH{FR32}, \LH{FR33}, \LH{FR34}, \LH{FR35}, \LH{FR37}, \LH{FR51}}
\begin{proposition}[Exact Certification Depends Only on the Decision Quotient Relation]\label{prop:depends-only-on-opt}
If two decision problems on the same state space induce the same decision-equivalence relation, then they induce the same exact-certification structure. In particular, they have the same sufficient coordinate sets, the same minimal sufficient coordinate sets, and the same relevant and irrelevant coordinates. Moreover, their quotient objects are canonically equivalent; in finite settings, those quotient objects therefore have the same cardinality.
\end{proposition}

Proposition~\ref{prop:depends-only-on-opt} is the canonical statement behind the frontier program. Exact certification is determined by the decision quotient relation itself. Equality of optimizer maps is only a sufficient route to this conclusion, because equal optimizer maps induce the same decision-equivalence relation.

\leanmetapending{\LH{FR48}}
\begin{corollary}[Sufficiency Is Relation Refinement]\label{cor:sufficiency-refinement}
A coordinate set $I$ is sufficient if and only if the coordinate-agreement relation induced by $I$ refines the decision quotient relation.
\end{corollary}

Corollary~\ref{cor:sufficiency-refinement} is the cleanest relational form of exact certification. It removes utility language entirely: sufficiency asks whether one state relation is contained in another.

\leanmetapending{\LH{FR49}, \LH{FR50}}
\begin{corollary}[Relevance Is Failure of Erased Sufficiency]\label{cor:relevance-erasure}
A coordinate $i$ is irrelevant if and only if deleting it leaves a sufficient coordinate set, namely $\mathrm{univ}\setminus\{i\}$. Equivalently, $i$ is relevant if and only if that erased set is not sufficient.
\end{corollary}

Corollary~\ref{cor:relevance-erasure} makes the role of individual coordinates canonical: relevance is exactly the obstruction to retaining sufficiency after coordinate erasure.

\leanmetapending{\LH{FR41}, \LH{FR42}, \LH{FR43}, \LH{FR44}}
\begin{corollary}[Certification Statistics Factor Through the Quotient Relation]\label{cor:statistics-factor}
Any statistic that is a function of the sufficient-set family and relevant-coordinate set is invariant under equality of the decision quotient relation. In particular, the number of sufficient coordinate sets, the number of minimal sufficient coordinate sets, and the number of relevant coordinates are quotient invariants.
\end{corollary}

Corollary~\ref{cor:statistics-factor} records the next logical consequence of Proposition~\ref{prop:depends-only-on-opt}. Once exact certification is identified with quotient data, every certification statistic built from sufficient sets or relevant coordinates factors automatically through that quotient data as well.

\leanmetapending{\LH{FR52}, \LH{FR53}}
\begin{corollary}[Constant-Optimizer Collapse Forces Trivial Certification]\label{cor:constant-opt-collapse}
If the optimizer set is constant across all states, then the empty coordinate set is sufficient and every coordinate is irrelevant.
\end{corollary}

Corollary~\ref{cor:constant-opt-collapse} isolates the first genuinely degenerate mechanism in the tractable basis. The certification problem disappears because there is no state dependence left to certify.

\leanmetapending{\LH{FR66}, \LH{FR67}}
\begin{proposition}[Explicit Enumeration Is Parameter-Dependent, Not Structural]\label{prop:enumeration-not-structural}
For every fixed exponent $k$, there exists a Boolean-cube family with state space size $2^n > n^k$. Thus explicit-state enumeration can outrun any fixed monomial bound in the ambient dimension parameter.
\end{proposition}

Proposition~\ref{prop:enumeration-not-structural} is weaker than a full polynomial lower-bound metatheorem, but it already captures the structural point needed here: tractability by brute-force state enumeration depends on a separate size parameter and does not arise from an optimizer-compatible mechanism like symmetry, low rank, or tree structure.

\leanmetapending{\LHrng{FR}{13}{16}}
\begin{proposition}[Positive Affine Utility Reparameterizations Are Invisible]\label{prop:affine-invisible}
Statewise positive affine transformations of utility leave exact certification unchanged. In particular, replacing
\[
U(a,s)
\qquad\text{by}\qquad
\alpha(s)+\beta(s)U(a,s)
\]
with $\beta(s)>0$ for every state $s$ does not change sufficient coordinate sets or relevance judgments.
\end{proposition}

Proposition~\ref{prop:affine-invisible} is another direct consequence of Proposition~\ref{prop:depends-only-on-opt}. It records a standard utility-theoretic invariance in the exact-certification language: only the induced decision quotient matters, not the particular positive affine scale used to encode utility.

\leanmetapending{\LH{FR54}, \LH{FR55}, \LH{FR56}, \LH{FR68}, \LH{FR69}, \LH{FR70}, \LH{FR71}, \LH{FR72}, \LH{FR73}}
\begin{proposition}[Duplicate Actions and Duplicate States Preserve Certification]\label{prop:duplicate-invariance}
Duplicating a single action without changing its utility profile preserves the decision quotient relation and hence preserves sufficiency and relevance. At the optimizer-set level, the only possible change is the addition of the duplicate action itself: original optimal actions remain optimal, and the duplicate is optimal exactly when the original action was. Duplicating a single state without changing its utility profile preserves the decision quotient up to the obvious quotient equivalence and also preserves sufficiency and relevance.
\end{proposition}

Proposition~\ref{prop:duplicate-invariance} gives two more theorem-forced closure laws. These are not modeling conventions; they are consequences of the fact that exact certification factors through quotient data rather than through accidental multiplicity in the presentation.

\leanmetapending{\LHrng{FR}{17}{20}}
\begin{proposition}[Relabeling Invariance Is Forced by Exact Certification]\label{prop:relabeling-forced}
Action relabeling preserves optimizer equivalence and exact sufficiency. State relabeling does so as well once the coordinate structure is transported along the state bijection. Consequently, any structural tractability frontier for exact relevance certification must be closed under these relabelings.
\end{proposition}

Proposition~\ref{prop:relabeling-forced} is now best read as a corollary of Proposition~\ref{prop:depends-only-on-opt}. Bijective relabelings do not change the optimizer map up to the obvious identification, so they cannot change exact certification.

The second necessary ingredient is a genuine expressivity restriction. Call a class \emph{kernel-universal} if, for every labeling map $\phi : S \to T$, it contains some decision problem whose optimizer quotient realizes exactly the kernel partition of $\phi$. By Theorem~\ref{thm:labeling-realizable}, kernel universality is an extremely strong property.

The canonical finite invariant here is the size of the optimizer quotient, equivalently the number of distinct optimal-action sets. The optimizer-quotient manuscript~\cite{simas2025optimizerquotient} identifies the quotient with the range of the optimizer map, so quotient size is not auxiliary bookkeeping; it is intrinsic.

\leanmetapending{\LH{FR24}, \LH{FR38}}
\begin{proposition}[Quotient Size Is Unbounded Under Realizability]\label{prop:quotient-size-unbounded}
For every $m \in \mathbb{N}$, there exists a finite decision problem whose optimizer quotient has size exactly $m$.
\end{proposition}

\begin{proof}
Take $m=k+1$ and realize the identity labeling on $\mathrm{Fin}(k+1)$. Proposition~\ref{prop:realized-quotient-range} identifies the quotient object with the range of that identity labeling, which has cardinality $k+1$. Equivalently, Proposition~\ref{prop:setoid-realizable} realizes the discrete equivalence relation on $\mathrm{Fin}(k+1)$.
\end{proof}

The point is simple: a meaningful frontier statement cannot range over classes that are simultaneously relabel-invariant and kernel-universal, because quotient shape alone then becomes too expressive to isolate the boundary.

\leanmetapending{\LH{FR21}, \LH{FR22}, \LH{FR23}}
\begin{proposition}[Invariance Alone Is Too Weak]\label{prop:invariance-too-weak}
The universal class of all decision problems satisfies action relabeling invariance and state relabeling invariance, but it is kernel-universal. Consequently, relabeling invariance by itself is far too weak to isolate a tractability frontier.
\end{proposition}

Proposition~\ref{prop:invariance-too-weak} separates two kinds of conditions that can otherwise be conflated.

\begin{enumerate}
\item \textbf{Benign invariance axioms} say the family should ignore arbitrary naming choices.
\item \textbf{Expressivity-limiting axioms} say the family should not realize arbitrary quotient structure wholesale.
\end{enumerate}

Any genuine optimizer-compatible frontier theorem will need both. Without decision-quotient invariance and its immediate corollaries such as relabel invariance and positive affine invariance, the statement is not structural. Without an expressivity restriction excluding kernel universality or something comparably strong, the realizability theorem and Proposition~\ref{prop:quotient-size-unbounded} make the class too large for quotient shape alone to carry the boundary.
